\section{O That Thou Wouldst Rend the Heavens}
\label{sec:opening}

The Church opens its year with a prayer that God would break the sky open.
Over the first reading of this Sunday the Latin \work{Ordo lectionum Missae},
on which the Lectionary rests, sets a single line of Isaiah as its title:
\latin{Utinam dirumperes caelos et descenderes!}, ``O that thou wouldst rend
the heavens, and wouldst come down'' (Is 63:19b; Douay 64:1). The people
who speak in the reading call the Lord their father, ``our redeemer, from
everlasting is thy name'', and ask him, ``Why hast thou made us to err, O
Lord, from thy ways?'' (Is 63:16b--17). In a verse the reading passes over,
they say that their enemies have trodden down his sanctuary (63:18). A few
lines after the cry, the Latin of the reading has them say something else:
\latin{Descendisti}, ``thou didst come down'' (Is 64:2b; Douay 64:3). Jerome,
who translated the prophet from the Hebrew, heard in that past tense a prayer
that had been granted.

{\small\noindent Scripture is quoted throughout in the Douay--Rheims version
as revised by Bishop Challoner, a public-domain study translation made from
the Latin Vulgate. It is not the text proclaimed at Mass: the English heard in
the celebration is that of the \work{Lectionary for Mass for Use in the
Dioceses of the United States of America}, which alone is the approved
wording. Isaiah is cited by the Lectionary's chapters and verses, which follow
the Hebrew division, with the Douay--Rheims number beside it where it
differs; psalms are numbered as the Lectionary numbers them, with the Vulgate
number, which the Douay--Rheims and the Latin books use, in parentheses.
Where Augustine on the psalms, Bellarmine or a Father who wrote in Greek is
quoted in English, the English is that of the translation named in the
References; other translations from Latin are the study's own.\par}
\medskip

The English proclaimed in the dioceses of the United States does not contain
those words: a congregation that hears the reading in English hears a people asking God to
come and does not hear them say that he came.

In the Gospel the one who has come speaks of his return. On the Mount of Olives,
after telling four of his disciples that no one knows the day or the hour,
Jesus commands them to stay awake for a master who has gone abroad and may
return at any watch of the night: ``And what I say to you, I say to all:
Watch'' (Mk 13:37). Between the prophet and the Gospel the psalm asks the
Shepherd of Israel to stir up his might and come, and Paul thanks God for a
church that lacks no gift as it waits for the revelation of Christ, because
``God is faithful'' (1~Cor 1:9). The Missal's prayers, the same in all three
years of the Lectionary's cycle, ask for the will to go out to meet Christ as
he comes, and the Preface gives thanks for his two comings, the first in
lowly flesh and the second in glory.

Three questions rise from these texts. The first is whether the prophet's cry
has been answered, and if so, how the Church can still pray it. The second is
why Christ hides the hour of his coming, and to whom his command is spoken
when most who hear it will die before the end. The third is on what ground a
people that has strayed can hope to be found blameless on the day.

\subsection{The texts of the Mass at a glance}

\begin{maptable}
Entrance Antiphon & \latin{Ad te levavi animam meam}; cf.\ Ps 25 (24):1--3,
an adaptation &
The speaker lifts up his soul to God, trusts in him, asks not to be shamed or
mocked by his enemies, and says that none who wait for God will be
confounded.\\
Collect & \latin{Da, quaesumus, omnipotens Deus, hanc tuis fidelibus
voluntatem} &
That God give his faithful the will to go out with just works to meet his
Christ as he comes, so that, set at his right hand, they may be counted worthy
of the kingdom of heaven.\\
First Reading & Is 63:16b--17, 19b; 64:2b--7 (Douay 63:16b--17; 64:1,
3--8) &
The people call the Lord their father and redeemer and ask why he has let
them stray; they cry to him to rend the heavens and come down; they confess
that he meets the just, that he was angry and they sinned, that their justice
is a soiled rag and no one calls on his name; they are clay, and he the
potter.\\
Responsorial Psalm & Ps 80 (79):2a and 2c, 3b, 15--16, 18--19; response
from v.~4 &
The Shepherd of Israel, enthroned on the cherubim, is asked to shine forth,
stir up his might and come; to turn, look down from heaven and visit the
vine; to lay his hand on the man of his right hand, the son of man; the
people promise not to turn away. The response asks God to turn his people and
show his face.\\
Second Reading & 1~Cor 1:3--9 &
Paul greets the Corinthians and thanks God for the grace given them in
Christ, in whom they lack no gift as they wait for his revelation; Christ will
confirm them to the end, blameless on his day; God, who called them into the
fellowship of his Son, is faithful.\\
Alleluia verse & Ps 85 (84):8 &
In the Douay--Rheims, ``Shew us, O Lord, thy mercy; and grant us thy
salvation.''\\
Gospel & Mk 13:33--37 &
Jesus commands watchfulness because the time is unknown; a man going abroad
leaves his servants authority and his porter the watch; the master may come
at any of the night's four watches; what he says to the disciples he says to
all: Watch.\\
Prayer over the Offerings & \latin{Suscipe, quaesumus, Domine, munera} &
That the Lord accept gifts offered from what he has given, and that the
devotion he grants in this passing time may become for those who offer it the
reward of eternal redemption.\\
Preface & Preface I of Advent, \latin{De duobus adventibus Christi} &
Thanks through Christ, who came first in the lowliness of flesh to open the
way of salvation, and who will come again in glory, when those who now keep
watch will receive openly what they were promised.\\
Communion Antiphon & \latin{Dominus dabit benignitatem}; Ps 85 (84):13 &
The psalm verse without its first word; in the Douay--Rheims, ``the Lord will
give goodness: and our earth shall yield her fruit''.\\
Prayer after Communion & \latin{Prosint nobis, quaesumus, Domine, frequentata
mysteria} &
That the mysteries received may profit those who receive them, since through
them God teaches them already, among passing things, to love what is
heavenly and hold to what lasts.\\
Solemn blessing (optional) & \latin{Omnipotens et misericors Deus, cuius
Unigeniti adventum} &
Three invocations for those who believe that the Son has come and await his
coming again: that God sanctify them, strengthen them in faith, hope and
charity, and reward them with eternal life when he comes in majesty.\\
\end{maptable}

\noindent The Gloria is omitted and the Creed is said. The Preface is
appointed by the Missal's rubric for this Mass; the solemn blessing may be
used in place of the simple blessing and is the one option among the texts
above. The Lectionary gives one form of each reading. A sung chant approved
for the purpose may replace either antiphon, the psalm or the Alleluia verse;
where one does, what is said below of that text's wording does not hold for
that celebration.
